% ===================================================================
%  Section III: Threat Model and Problem Formulation
% ===================================================================

\subsection{System Roles}

The system defines five roles. A \emph{dataset collector/importer} creates
replay data and metadata through either instrumented environment interaction or
public benchmark import. A \emph{prover} opens committed transitions and
model/checkpoint witnesses and produces relation proofs or semantic-verifier
evidence. A \emph{verifier} checks public statements, proof public values, and
artifact hashes. An \emph{auditor/report generator} checks source artifacts and
emits benchmark tables. A \emph{benchmark runner} records proof cost and
tamper-rejection evidence without changing the relations.

\subsection{Assets, Public Inputs, and Private Witnesses}

The protected assets are replay dataset contents, dataset commitments, manifest
and audit reports, model weights, checkpoint hashes, private witness values,
proof artifacts, and benchmark tables. \Cref{tab:public-private} summarizes the
public/private classification.

\begin{table}[!t]
\centering
\small
\caption{Public inputs vs.\ private witnesses in the proof system.}
\label{tab:public-private}
\begin{tabular}{p{3.8cm}p{3.8cm}}
\toprule
\textbf{Public inputs} & \textbf{Private witnesses} \\
\midrule
Dataset root / Merkle root &
Transition field values \\
Manifest hash, audit-report hash &
Merkle sibling paths \\
Raw trajectory hash &
Model weights and intermediates \\
Config hash &
Gradient tensors \\
Checkpoint hashes (in/out) &
Fixed-point intermediate values \\
Relation identifier and outputs &
Minibatch indices (unless public) \\
Proof public values &
--- \\
\bottomrule
\end{tabular}
\end{table}

The private witness boundary is relation-specific. Transition values may be
private witnesses to membership and training relations, but they can also be
public or committed elsewhere. Merkle sibling paths are witnesses unless a
report discloses them. Model weights and intermediate tensors are private when
only checkpoint hashes are public.

\subsection{Adversary Model}\label{subsec:adversary}

The adversary is a computationally bounded malicious prover who controls the
training process and seeks to produce a convincing but dishonest proof. The
adversary may attempt any combination of the following attacks:

\begin{enumerate}[leftmargin=1.8em, label=\textbf{A\arabic*}.]
  \item \textbf{Data tampering}: alter a transition reward, action, done flag,
  or next state; use a transition outside the committed replay set.
  \item \textbf{Commitment tampering}: tamper with a Merkle path, leaf index,
  leaf hash, sibling hash, or root.
  \item \textbf{Sampling tampering}: change a deterministic minibatch index or
  duplicate an index within a batch.
  \item \textbf{Computation tampering}: alter Q values, TD targets, losses,
  SmoothL1 derivatives, gradients, SGD deltas, or weight updates.
  \item \textbf{Trace tampering}: skip or reorder checkpoint steps; break
  checkpoint hash chains; tamper with target-network synchronization timing.
  \item \textbf{Aggregation tampering}: reorder aggregation chunks; swap child
  proof manifests; alter child public-input hashes; present inconsistent
  dataset or configuration commitments.
  \item \textbf{Verification tampering}: verify a proof against different
  public inputs than those committed.
\end{enumerate}

The verifier accepts only if the implemented relation and public-input binding
checks pass for the specific relation being verified.

\subsection{Reward and Dataset Correctness Boundary}

For self-collected datasets, the audit pipeline replays the environment and
checks reward, next-state, termination, truncation, and collection-log evidence
\emph{before} commitment. After commitment, proofs and verifiers check
consistency with the committed and audited data. The system does not establish
reward correctness for arbitrary external datasets beyond the import and
source-integrity checks that were performed. Public Minari~\cite{minari2023}
and D4RL~\cite{fu2020d4rl} benchmark rows are source-integrity commitments,
not honest-collection proofs.

\subsection{Security Goals and Non-Goals}

\noindent\textbf{Goals.} The system aims to provide:
\begin{itemize}[leftmargin=1.5em]
  \item Soundness of replay membership relative to a committed dataset root.
  \item Correctness of fixed-point TD/Bellman computation for opened
  transitions.
  \item Correctness of MLP forward pass, backpropagation, and SGD update for
  canonical tiny vectors.
  \item Integrity of checkpoint chains and target-network synchronization
  within training fragments.
  \item Consistency of proof-manifest chunk-chain aggregation.
  \item Zero-knowledge privacy for witness data not intentionally published.
\end{itemize}

\noindent\textbf{Non-goals.} The system explicitly does \emph{not} claim:
\begin{itemize}[leftmargin=1.5em]
  \item Full DQN training proof from initialization to final checkpoint.
  \item Adam optimizer or advanced optimizer-state proofs.
  \item All-replay-batches proof or model-selection proof.
  \item True recursive child-proof verification inside the aggregate SP1 guest.
  \item Honest public dataset collection proof.
  \item Policy optimality or convergence guarantees.
\end{itemize}

\subsection{Security Claims Matrix}

\Cref{tab:security-claims} maps each threat category to the corresponding
mitigation mechanism, formal theorem, tamper-rejection evidence
(\Cref{sec:experiments}), and known limitation.

\begin{table*}[!t]
\centering
\small
\caption{Security claims matrix: threats, mitigations, theorem coverage, and residual limitations.}
\label{tab:security-claims}
\begin{tabular}{p{2.8cm} p{3.8cm} p{1.8cm} p{2.8cm} p{3.6cm}}
\toprule
\textbf{Threat} & \textbf{Mitigation} & \textbf{Theorem} &
\textbf{Table~\ref{tab:tamper-summary} category} & \textbf{Residual limitation} \\
\midrule
Use transition outside replay set &
Merkle leaf/path/root recomputation &
Thm.~\ref{thm:replay-membership} &
Merkle path, dataset root &
No honest collection proof for public data \\
\midrule
Alter reward, action, done, or next state &
Dataset audit + TD relation recomputation &
Thms.~\ref{thm:dataset-commitment}, \ref{thm:bellman-target} &
reward, action, done, next\_state &
Reward audit only where collection artifacts support it \\
\midrule
Falsify Q value, TD target, or loss &
Fixed-point Bellman and SmoothL1 recomputation &
Thm.~\ref{thm:bellman-target} &
Q value, TD target, loss &
No policy optimality claim \\
\midrule
Falsify gradient or update &
Fixed-point backpropagation and SGD checks &
Thm.~\ref{thm:update-correctness} &
gradient, checkpoint hash &
Batch-size-1 tiny MLP only \\
\midrule
Break checkpoint chain or target sync &
Per-step checkpoint + target checkpoint links &
Thms.~\ref{thm:checkpoint-chain}, \ref{thm:training-fragment} &
checkpoint hash, target sync &
No full training-run proof \\
\midrule
Reorder chunks / swap proof manifests &
Manifest hashes + aggregate roots + checkpoint links &
Thm.~\ref{thm:manifest-aggregation} &
aggregation, proof public input &
Proof-manifest chain only \\
\midrule
Verify under different public inputs &
Public-input binding and proof public values &
Thm.~\ref{thm:privacy-boundary} &
proof public input &
Published metadata remains public \\
\bottomrule
\end{tabular}
\end{table*}
